\documentclass[11pt]{article}

\usepackage{amsmath,amssymb,graphicx}
\usepackage[margin=1in]{geometry}
\usepackage{booktabs}

\title{Relational Feedback Dynamics Probe on a Frozen Kernel-Induced Cochain Operator Architecture}
\author{Tomislav Rupi\'c \\ Independent Researcher}
\date{March 2026}

\begin{document}

\maketitle

\begin{abstract}
Stages 10--15 established a frozen pre-geometric ladder on a kernel-induced cochain operator architecture: morphology, collective interaction, coarse envelope structure, scale-coupled fragility, falsified weak binding channels, and a first relational-ordering probe. Stage 15 found that metric-like ordering can appear, but only in a regime-dependent way: the cleanest ordering was observed in simple propagation corridors and phase-ordered diffuse configurations, while the strongest composite anchor remained weak or metric-violating. Stage 16 asks the next falsification-first question: can such derived relational ordering act back on dynamics as a weak organizing influence? A narrow nine-run pilot is executed on three representative cases with three conditions per case: null control, relational feedback strength $\varepsilon=0.01$, and relational feedback strength $\varepsilon=0.02$. The feedback term is constructed only from frozen Stage 15 relational drivers and leaves the underlying kernel, projector, and base operator unchanged. The result is cleanly negative. All nine base runs return ``no relational dynamic effect''; no interaction label changes, no coarse label changes, no composite lifetime gain appears, and no case earns promotion to the refinement ladder. The bounded conclusion is that, at the tested strengths and with the present relational constructions, Stage 15 ordering does not yet condense into an effective organizing dynamics.
\end{abstract}

\section{Introduction}

The program through Stage 15 is cumulative and deliberately restrictive. Stage 10 built a morphology instrument for frozen packet transport. Stage 11 extended that instrument to collective interaction patterns. Stage 12 coarse-grained multi-packet fields into envelope and basin structure. Stage 13 showed that coarse organization sharpens under refinement while composite interaction labels remain fragile. Stage 14 then tested three weak structural branches and found no effective binding channel at the tested strengths. Stage 15 changed viewpoint and asked whether frozen interaction data already support a stable relational ordering that behaves metric-like under refinement. The resulting signal was mixed but informative: some localized metric-like ordering appeared, but not on the strongest composite anchor.

Stage 16 asks the next question in the same falsification-first style. If relational ordering is beginning to condense into a more structural layer, then feeding a weak version of that ordering back into the dynamics should selectively stabilize the cases that already showed cleaner ordering in Stage 15. If no such selective stabilization appears, then the relational ordering remains diagnostic bookkeeping rather than an effective organizing influence.

\section{Frozen Context and Pilot Design}

All kernels, projectors, operator constructions, and packet families remain frozen. No new interaction term is added at the level of the kernel or cochain operator. The projected transverse sector and periodic boundary condition are retained, together with the packet family inherited from the strongest coherent transverse seed: bandwidth $w=0.1$, cosine carrier, central wave number $k=1.0$, and amplitude scale $0.5$.

The Stage 16 pilot uses exactly three cases:
\begin{itemize}
\item phase-ordered symmetric triad,
\item counter-propagating corridor,
\item clustered composite anchor.
\end{itemize}

Each case is paired with a single relational driver inherited from Stage 15:
\begin{itemize}
\item phase-correlation distance for the symmetric triad,
\item transport-delay distance for the corridor,
\item persistence distance for the composite anchor.
\end{itemize}

Each case is run under three conditions:
\begin{itemize}
\item null control,
\item relational feedback with $\varepsilon=0.01$,
\item relational feedback with $\varepsilon=0.02$.
\end{itemize}

This gives a total of nine base runs. The Stage 16 design allows a $(12 \rightarrow 24)$ refinement follow-up only if a base-resolution feedback run shows measurable stabilization over null without global artificial binding. No run satisfied that gate.

\section{Relational Feedback Construction}

The feedback is built only from frozen Stage 15 relational structures. For packet $i$, the schematic form is
\[
F_{\mathrm{rel}}(i) = \varepsilon \sum_j W_{\mathrm{rel}}(i,j)\,(x_j - x_i),
\]
where $W_{\mathrm{rel}}$ is row-normalized and constructed from the selected Stage 15 driver for the case. The implementation is bounded, local in the induced relational-graph sense, and applied as a weak directional modulation of the packet component rather than as a modification of the frozen base operator.

Three drivers are used:
\begin{itemize}
\item persistence-overlap weights,
\item phase-lock weights,
\item transport-response weights.
\end{itemize}

No claim is made that this is a physical force law or a metric dynamics. The purpose is only to test whether relational bookkeeping can begin to act back on morphology in a selective and measurable way.

\section{Measurement Set}

The pilot reuses the collective and coarse observables frozen in earlier stages:
\begin{itemize}
\item composite lifetime,
\item pair-distance evolution,
\item overlap persistence,
\item dominant basin area fraction,
\item basin-count stability,
\item interaction label,
\item coarse label.
\end{itemize}

It adds three Stage 16 summary indicators:
\begin{itemize}
\item relational stabilization score,
\item morphology-drift reduction,
\item feedback-induced regime change.
\end{itemize}

Each run is assigned one of four descriptive output labels:
\begin{itemize}
\item no relational dynamic effect,
\item weak stabilization,
\item selective regime stabilization,
\item global artificial binding.
\end{itemize}

\section{Pilot Outcome}

The Stage 16 result is uniform:
\begin{itemize}
\item no relational dynamic effect: 9,
\item weak stabilization: 0,
\item selective regime stabilization: 0,
\item global artificial binding: 0.
\end{itemize}

No case is promoted to the refinement ladder.

\subsection{Phase-ordered symmetric triad}

The phase-ordered triad is the strongest non-degenerate Stage 15 phase-ordering carrier. Nevertheless, under both tested feedback strengths it remains unchanged at the reported resolution. All three conditions return the same interaction label, \emph{transient binding regime}, and the same coarse label, \emph{diffuse coarse field regime}. Composite lifetime, binding persistence, and coarse persistence remain unchanged relative to null at the recorded precision.

\subsection{Counter-propagating corridor}

The counter-propagating corridor carried the cleanest localized Stage 15 ordering signal, including the only stable transport-delay ordering in the pilot. Even here, the Stage 16 feedback produces no measurable response. All three conditions remain in the \emph{transient binding regime} with \emph{diffuse coarse field regime} at the coarse layer. No lifetime gain or basin-stability gain is detected.

\subsection{Clustered composite anchor}

The clustered composite anchor is retained precisely because it had the strongest collective composite morphology while remaining weak or metric-violating in Stage 15. In Stage 16 it stays unchanged under both feedback strengths. All three conditions remain in the \emph{metastable composite regime} with \emph{diffuse coarse field regime} at the coarse layer. No increase in composite lifetime or binding persistence is observed.

\section{Interpretation of the Null Result}

The null result is structurally useful. Stage 15 showed that some metric-like ordering can be extracted from the frozen data, but Stage 16 shows that this ordering does not yet act back on the dynamics in a measurable way at the tested strengths. The key point is selective failure: the feedback does not stabilize the cases that already carried cleaner ordering, but it also does not create spurious global binding across all cases.

This means the present relational constructions remain diagnostic rather than dynamical. They organize interpretation of the frozen data, but they do not yet function as an effective organizing layer for morphology.

The result also sharpens the overall program boundary. By Stage 14 the weak binding channel had already been falsified under several weak extensions. Stage 16 now adds a second negative result at a different conceptual layer: even when weak relational order is explicitly fed back into the dynamics, no measurable stabilization appears in the tested regimes.

\section{Interpretation Boundary}

Stage 16 does not support:
\begin{itemize}
\item spacetime claims,
\item physical metric claims,
\item causal-structure claims,
\item effective-force claims.
\end{itemize}

The only bounded claims supported here are:
\begin{itemize}
\item Stage 15 relational ordering can be constructed diagnostically,
\item at $\varepsilon=0.01$ and $0.02$ that ordering does not measurably reorganize collective morphology,
\item no selective stabilization appears in the already ordered regimes,
\item no global artificial binding is induced either.
\end{itemize}

\section{Conclusion}

Stage 16 is the first direct probe of whether derived relational ordering can feed back into dynamics as a weak organizing influence on the frozen architecture. The answer, in this first executable pilot, is negative. Across nine base runs spanning a phase-ordered triad, a transport-dominated corridor, and a clustered composite anchor, the tested relational feedback produces no measurable morphology response.

This is a useful constraint rather than a dead end. The program now distinguishes clearly among three layers:
\begin{itemize}
\item morphology and interaction structure can exist,
\item relational ordering can sometimes be read out diagnostically,
\item that ordering does not yet act as an effective dynamic organizer under the tested weak feedback channel.
\end{itemize}

The next sensible step is therefore not stronger geometry language. It is a more careful refinement of relational transport structure and ordering resilience under perturbation and scale stress.

\end{document}
